\section{Future work}

\subsection{Timer support}

\subsection{External interrupts}

\subsection{Safeguarding system calls}
With the current setup, when a new domain is created, it must be created with a
task that is in privileged, if that domain should have more than a single task
running. The reason for this is two fold. First, to create a task in
FreeRTOS-MPU one must be in privileged mode. Secondly, to add a task to a
domain with the new API presented in Section~\ref{sec:api}, one must either be
inside the domain itself or be creating a new domain\footnote{Given the task
  runs in privileged mode, it could technically overwrite the memory locations
forcefully, and not make use of the api at all.}. The requirement of a single
privileged task mainly comes from the FreeRTOS-MPU version, and to circumvent
this issue would require significant changes from FreeRTOS itself. The larger
issue in regards to domain creation is, that creating a new domain could delay
scheduling. Specifically, many of the privileged functions enter a "critical"
section within FreeRTOS. When entering such a critical section, on the RISC-V
MPU version interrupts are disabled, including the machine timer interrupts
which drives the domain scheduling. This poses a significant burden on the
privileged tasks within each domain, as they can inadvertently create timing
side-channels. The issue is not only with the domain creation, but any call to
functions that enter the critical sections. 

This issue poses a few challenges, and there are two primary solutions, which
could be considered for future development. The first is taking inspiration
from \cite{ge2019}. Here, they clone the shared kernel data into memory
provided by user level tasks within each domain. The current implementation
created within FreeRTOS does not do this, but many of the preexisting features
of the within domain priority scheduler could be fully partitioned. For this to
function, one would have to introduce two different kinds of critical sections.
There is the cross-domain critical section, for which we would still have to
disable interrupts when it occurs, and there is the within domain critical
section. When entering a within domain critical section, then we would block
any priority scheduling, such that the priority based scheduler is not left in
an invalid state, but if there is a domain level interrupt, then we would save
the context of the current execution and switch to another domain. When the old
domain runs again, it would still be within the domain critical section at the
same context point, allowing for a valid state of the within domain scheduler,
while allowing the domain level scheduler to schedule other domains without
delays. This context saving and restoration would have to be taken into account
for the worst case analysis of the domain switching, if it can not be
guaranteed in constant time, but should propagate well to the majority of
FreeRTOS' standard implementation. 

The next issue would be the cross-domain critical section. This part of the
kernel is shared between all, and as such we can not simply "ignore" the
critical section as before. Instead we would have to guarantee that any call to
cross domain critical functions would execute within the current domain. A
solution to this problem would be to guarantee, that any system call with such
a cross-domain critical section runs within a single time slot, or allows for
preemption within a time slice. Then, when entering the cross-domain critical
section, the task would have to idle until it reaches the start of a time slice
and then execute the system call. This would mean tat execution is guaranteed
to finish before the scheduler would attempt to schedule the next domain. 

For this to function correctly one would need proper hardware to find the
running time of all such system calls, and make sure that the execution time is
bounded by the time slice. It would also mean, that there has to be a hard
requirement on the lowest possible value of the time slice duration. However,
with these in mind truly isolated domains should be possible.

%TODO: Possible this is a channel in S3k

